% from ega4-2-fr.tex
\oldpage[IV-2]{49}

\begin{corollary}[3.4.2]
\label{IV.2.3.4.2-fr}
Sous les hypothèses générales de (3.4.1), supposons que $Z$ ne soit pas égal
à un des $S_i$ et que
$\operatorname{long}((\mathcal F/t\mathcal F)_z)=1$. Alors il existe un seul
des $S_i$ contenant $Z$, et pour cette valeur de $i$, on a
$\operatorname{long}(\mathcal F_{s_i})=1$~; en outre
$\mathcal O_{S_i,z}$ est un anneau de valuation discrète dont $t_z$ est une
uniformisante, et $t_z$ est $\mathcal F_z$-régulier.
\end{corollary}

\begin{proof}
Cela résulte de (3.4.1), les deux membres de (3.4.1.1) étant alors égaux.
\end{proof}

\begin{proposition}[3.4.3]
\label{IV.2.3.4.3-fr}
Soient $X$ un préschéma localement noethérien, $t$ une section de
$\mathcal O_X$ au-dessus de $X$, $Y$ le sous-préschéma fermé de $X$ défini par
l'Idéal $t\mathcal O_X$ de $\mathcal O_X$. Soient $\mathcal F$ un
$\mathcal O_X$-Module cohérent, $T$ un cycle premier associé à $\mathcal F$,
$T'$ une composante irréductible de $T\cap Y$, $x$ le point générique de
$T'$. Supposons que $t_x$ soit $\mathcal F_x$-régulier~; alors on a
$x\in\operatorname{Ass}(\mathcal F/t\mathcal F)$.
\end{proposition}

\begin{proof}
Comme dans la démonstration de (3.4.1), on peut se ramener au cas où
$X=\operatorname{Spec}(\mathcal O_x)$~; la proposition est alors (compte tenu
de (3.1.2)) conséquence de $(0,16.4.6.3)$.
\end{proof}

\begin{proposition}[3.4.4]
\label{IV.2.3.4.4-fr}
Soient $X$ un préschéma localement noethérien, $t$ une section de
$\mathcal O_X$ au-dessus de $X$, $Y$ le sous-préschéma fermé de $X$ défini par
l'Idéal $t\mathcal O_X$ de $\mathcal O_X$. Soient $\mathcal F$ un
$\mathcal O_X$-Module cohérent, $(S_i)$ la famille des composantes
irréductibles de $\operatorname{Supp}(\mathcal F)$. Soit $y$ un point de $Y$
tel que $t_y$ soit $\mathcal F_y$-régulier et qu'aucun des cycles premiers
associés immergés de $\mathcal F/t\mathcal F$ ne contienne $y$. Alors les
composantes irréductibles de $\operatorname{Supp}(\mathcal F/t\mathcal F)$ qui
contiennent $y$ sont exactement les composantes irréductibles des $S_i\cap Y$
qui contiennent $y$, et les cycles premiers associés à $\mathcal F$ qui
contiennent $y$ sont non immergés.
\end{proposition}

\begin{proof}
Prouvons d'abord la dernière assertion. Soient $T\supset T_1$ deux cycles
premiers associés à $\mathcal F$ contenant $y$~; si $x$ est point générique
d'une composante irréductible de $T\cap Y$ contenant $y$, $x$ est
générisation de $y$, donc contenu dans tout voisinage de $y$, et l'hypothèse
que $t_y$ est $\mathcal F_y$-régulier entraîne que $t_x$ est
$\mathcal F_x$-régulier $(0,15.2.4)$, donc, en vertu de (3.4.3), on a
$x\in\operatorname{Ass}(\mathcal F/t\mathcal F)$. Soit $x_1$ le point
générique d'une composante irréductible de $T_1\cap Y$ contenant $y$, et soit
$x$ le point générique d'une composante irréductible de $T\cap Y$ contenant
$x_1$~; il résulte de ce qui précède que $x_1$ et $x$ appartiennent tous deux
à $\operatorname{Ass}(\mathcal F/t\mathcal F)$, et comme
$x_1\in\overline{\{x\}}$, l'hypothèse entraîne que $x_1=x$. Désignons encore
par $T$ et $T_1$ les sous-préschémas fermés intègres de $X$ ayant $T$ et
$T_1$ comme espaces sous-jacents respectifs, et posons
$A=\mathcal O_{T,x}$, $A_1=\mathcal O_{T_1,x}$~; on a donc
$A_1=A/\mathfrak p$, où $\mathfrak p$ est un idéal premier de $A$. Par
définition de $x$ et $x_1$, $A/tA$ et $A_1/tA_1$ sont des anneaux artiniens~;
d'autre part, on a vu plus haut que $t_x$ est $\mathcal F_x$-régulier, donc
$x$ ne peut appartenir à $\operatorname{Ass}(\mathcal F)$, et par suite $A_1$
n'est pas artinien. On a donc $\dim A=\dim A_1=1$~; mais cela entraîne
$\mathfrak p=0$ et $A=A_1$ $(0,16.1.2.2)$~; comme
$\operatorname{Spec}(A)$ et $\operatorname{Spec}(A_1)$ sont respectivement
denses dans $T$ et $T_1$, on a bien $T=T_1$.

Les $S_i$ qui contiennent $y$ sont donc tous les cycles premiers associés à
$\mathcal F$ contenant $y$~; si $x$ est le point générique d'une composante
irréductible de $S_i\cap Y$ contenant $y$, on déduit encore de $(0,15.2.4)$
que $t_x$ est $\mathcal F_x$-régulier, donc, par (3.4.3), que
$x\in\operatorname{Ass}(\mathcal F/t\mathcal F)$~; cela démontre la première
assertion de (3.4.4).
\end{proof}

\begin{proposition}[3.4.5]
\label{IV.2.3.4.5-fr}
Soient $X$ un préschéma localement noethérien, $t$ une section de
$\mathcal O_X$ au-dessus de $X$, $Y$ le sous-préschéma fermé de $X$ défini par
l'Idéal $t\mathcal O_X$ de $\mathcal O_X$. Soient $\mathcal F$

